\documentclass[11pt]{article}
\usepackage[margin=18mm,headheight=22pt]{geometry}
\usepackage{amsmath}
\usepackage{fontspec}
\usepackage{unicode-math}
\setmainfont{NewCM10}[Extension=.otf,UprightFont=*-Book,BoldFont=*-Bold,ItalicFont=*-BookItalic,BoldItalicFont=*-BoldItalic]
\setsansfont{NewCMSans10}[Extension=.otf,UprightFont=*-Book,BoldFont=*-Bold,ItalicFont=*-BookOblique,BoldItalicFont=*-BoldOblique]
\setmathfont{NewCMMath-Book.otf}
\usepackage{xcolor}
\definecolor{FAUblue}{HTML}{04316A}
\definecolor{Paperblue}{HTML}{EEF2F7}
\definecolor{EEIgold}{HTML}{CC9C42}
\usepackage{etoolbox,enumitem,booktabs,tabularx,array,needspace,microtype}
\usepackage[most]{tcolorbox}
\usepackage{titlesec,fancyhdr}
\usepackage[hidelinks]{hyperref}
\usepackage{draftwatermark}
\SetWatermarkText{Made by Thanate Amornjiraporn (Focus)}
\SetWatermarkColor{black!7}\SetWatermarkFontSize{30pt}\SetWatermarkAngle{45}
\setlength{\parindent}{0pt}
\setlength{\parskip}{5pt}
\setlist{itemsep=2pt,topsep=3pt,leftmargin=1.6em,beginpenalty=10000,endpenalty=10000}
\raggedbottom
\clubpenalty=10000 \widowpenalty=10000 \displaywidowpenalty=10000 \interlinepenalty=10000
\setcounter{secnumdepth}{0}
\setcounter{tocdepth}{1}
\pagestyle{fancy}\fancyhf{}
\fancyhead[L]{\small\sffamily\color{FAUblue} From Fundamentals to AI Apocalypse · APDR · Blackboard Exercise 5}
\fancyhead[R]{\small\sffamily Blackboard Exercise 5: Worked solutions}
\fancyfoot[C]{\small\thepage}
\titleformat{\section}{\Large\bfseries\sffamily\color{FAUblue}}{}{0pt}{}
\titleformat{\subsection}{\large\bfseries\sffamily\color{FAUblue}}{}{0pt}{}
\titlespacing*{\section}{0pt}{16pt}{6pt}
\preto\section{\Needspace{8\baselineskip}}
\preto\subsection{\Needspace{6\baselineskip}}
\newtcolorbox{equationcard}[1]{colback=Paperblue,colframe=Paperblue,boxrule=0pt,arc=0pt,left=10pt,right=10pt,top=6pt,bottom=6pt,title={#1},fonttitle=\small\bfseries\sffamily,coltitle=FAUblue,before skip=8pt,after skip=8pt}
\newtcolorbox{workedbox}[1]{colback=white,colframe=EEIgold,boxrule=0pt,leftrule=3pt,arc=0pt,left=10pt,top=4pt,bottom=4pt,title={#1},fonttitle=\small\bfseries\sffamily,coltitle=FAUblue,attach title to upper={\par},before skip=8pt,after skip=8pt}
\newtcolorbox{recallbox}[1]{colback=white,colframe=FAUblue,boxrule=0pt,leftrule=3pt,arc=0pt,left=10pt,top=4pt,bottom=4pt,title={#1},fonttitle=\small\bfseries\sffamily,coltitle=FAUblue,attach title to upper={\par},before skip=8pt,after skip=8pt}
\usepackage{environ}
\newsavebox\measurebox
% Typeset the section once invisibly to measure it; start a new page if it will not fit in what is left.
\makeatletter
\newif\ifkeep@measuring
\NewEnviron{keepsection}{%
  \ifkeep@measuring \BODY \else
    \edef\keep@sec{\the\value{section}}\edef\keep@sub{\the\value{subsection}}%
    \setbox\measurebox\vbox{\hsize\linewidth\keep@measuringtrue\def\Needspace##1{}\BODY}%
    \setcounter{section}{\keep@sec}\setcounter{subsection}{\keep@sub}%
    \dimen2=\ht\measurebox \advance\dimen2\dp\measurebox
    \dimen@=\textheight
    \ifdim\pagegoal<\maxdimen
      \dimen@=\pagegoal \advance\dimen@-\pagetotal
      \ifdim\dimen2>\dimen@ \clearpage \dimen@=\textheight \fi
    \fi
    % if it fits, nothing inside may force an early break
    \ifdim\dimen2>\dimen@ \BODY \else \begingroup\def\Needspace##1{}\BODY\endgroup \fi
  \fi}
\makeatother
\newtcolorbox{notebox}{colback=Paperblue,colframe=Paperblue,boxrule=0pt,arc=0pt,left=8pt,right=8pt,top=4pt,bottom=4pt,before skip=6pt,after skip=6pt}
\newtcolorbox{plainbox}{colback=white,colframe=FAUblue,boxrule=0pt,leftrule=3pt,arc=0pt,left=10pt,right=8pt,top=4pt,bottom=4pt,before skip=6pt,after skip=8pt}
\usetikzlibrary{arrows.meta}
\title{\sffamily\bfseries\color{FAUblue}Blackboard Exercise 5: Worked solutions}
\author{Algorithms, Programming, and Data Representation \ensuremath{\cdot} Module 92313 \ensuremath{\cdot} Blackboard Exercise 5\\[2pt]\small From Fundamentals to AI Apocalypse · personal study guide}
\date{}
\begin{document}
\maketitle
\begingroup\footnotesize\setlength{\parskip}{0pt}\tableofcontents\endgroup
\clearpage

\Needspace{8\baselineskip}
{\large Build the recursive pattern from its base case and smaller calls, trace the Sierpinski counts, and use the call stack to move into and out of nested dream levels.}


\begin{keepsection}
\Needspace{5\baselineskip}
\section{Exercise 5.1: Recursion}

\textbf{Connected topic:} The recursive pattern (§8.1) \ensuremath{\cdot} Tracing calls (§8.2)

\begin{keepsection}
\Needspace{7\baselineskip}
\subsection{5.1.1 Define recursion}

\Needspace{4\baselineskip}
Try first: Define recursion in fewer than three sentences, including how a recursive process stops.

\nopagebreak[4]{\small\itshape Answers: see “Exercise 5.1: Recursion” under Answers at the end.}
\end{keepsection}

\begin{keepsection}
\Needspace{7\baselineskip}
\subsection{5.1.2 Complete the factorial function}

\Needspace{4\baselineskip}
Try first: Fill the blanks with a base-case comparison, its return value, and the decrement in the recursive call.

\nopagebreak[4]{\small\itshape Answers: see “Exercise 5.1: Recursion” under Answers at the end.}
\end{keepsection}
\end{keepsection}

\begin{keepsection}
\Needspace{5\baselineskip}
\section{Exercise 5.2: Sierpinski triangle}

\textbf{Connected topic:} The recursive pattern (§8.1) \ensuremath{\cdot} Tracing calls (§8.2)

\begin{keepsection}
\Needspace{7\baselineskip}
\subsection{5.2.1 Continue the triangle pattern}

\Needspace{4\baselineskip}
Try first: Continue the given stages for \textit{n} = 2 and \textit{n} = 3 while leaving the empty central triangles blank.

\nopagebreak[4]{\small\itshape Answers: see “Exercise 5.2: Sierpinski triangle” under Answers at the end.}
\end{keepsection}

\begin{keepsection}
\Needspace{7\baselineskip}
\subsection{5.2.2 Write a recursive triangle counter}

\Needspace{4\baselineskip}
Try first: Turn the repeated three-way growth into a recurrence with a base case for the starting triangle.

\nopagebreak[4]{\small\itshape Answers: see “Exercise 5.2: Sierpinski triangle” under Answers at the end.}
\end{keepsection}

\begin{keepsection}
\Needspace{6\baselineskip}
\subsection{5.2.3 Count the triangles at step 6}

\Needspace{3\baselineskip}
Try first: Use the recurrence table to find the black-triangle count at \textit{n} = 6.

\nopagebreak[4]{\small\itshape Answers: see “Exercise 5.2: Sierpinski triangle” under Answers at the end.}
\end{keepsection}
\end{keepsection}

\begin{keepsection}
\Needspace{5\baselineskip}
\section{Exercise 5.3: Recursion application}

\textbf{Connected topic:} The recursive pattern (§8.1) \ensuremath{\cdot} Tracing calls (§8.2)

\begin{keepsection}
\Needspace{6\baselineskip}
\subsection{5.3.1 Recursion and repetition}

\Needspace{3\baselineskip}
Try first: State the difference between recursion and repetition in one sentence.

\nopagebreak[4]{\small\itshape Answers: see “Exercise 5.3: Recursion application” under Answers at the end.}
\end{keepsection}

\begin{keepsection}
\Needspace{7\baselineskip}
\subsection{5.3.2 Write the recursive \texttt{dream} function}

\Needspace{4\baselineskip}
Try first: Make the dream level progress to the maximum, do the job once, unwind one level at a time, and check the totem after waking.

\nopagebreak[4]{\small\itshape Answers: see “Exercise 5.3: Recursion application” under Answers at the end.}
\end{keepsection}
\end{keepsection}


\clearpage
\section{Answers}
\begin{keepsection}\subsection{Exercise 5.1: Recursion}

\begin{workedbox}{5.1.1 Define recursion \ensuremath{\cdot} worked solution}
\Needspace{6\baselineskip}
\begin{enumerate}
\item \textbf{Recursion} solves a problem by having a function call itself on a simpler instance of the same problem; a base case answers directly, and each recursive case must make progress toward it.
\item \textbf{Check:} the definition includes both the self-call and the stopping case, so it rules out a call that repeats forever without getting simpler.
\end{enumerate}
\end{workedbox}

\end{keepsection}
\begin{keepsection}\subsection{Exercise 5.1: Recursion}

\begin{workedbox}{5.1.2 Complete the factorial function \ensuremath{\cdot} worked solution}
\Needspace{13\baselineskip}
\begin{enumerate}
\item Use the course's positive-integer convention for this version: \texttt{factorial(1)} is the base case, and for \textit{x} \ensuremath{>} 1 the recurrence is \textit{x}! = \textit{x} \ensuremath{\cdot} (\textit{x} \ensuremath{-} 1)!.
\item \Needspace{3\baselineskip}
Fill the blanks with \texttt{1}, \texttt{1}, and \texttt{1}:
\begin{tcolorbox}[colback=Paperblue,colframe=Paperblue,boxrule=0pt,arc=0pt,left=7pt,right=7pt,top=5pt,bottom=5pt]\small\sloppy
\noindent\hspace*{0em}\texttt{def factorial(x):}\par
\noindent\hspace*{1em}\texttt{if x == 1:}\par
\noindent\hspace*{3em}\texttt{return 1}\par
\noindent\hspace*{1em}\texttt{return x * factorial(x - 1)}\par
\end{tcolorbox}
\item For \texttt{factorial(5)}, calls descend to \texttt{factorial(1)}; then the waiting calls return 1, 2, 6, 24, and 120.
\item The sheet does not specify whether zero is an allowed input. If the function must also accept zero, use \texttt{x == 0} as the base case and return 1, which follows the convention 0! = 1.
\item \textbf{Check:} the argument decreases by one on each call, so every positive integer input reaches 1; the value then multiplies back up.
\end{enumerate}
\end{workedbox}

\end{keepsection}
\begin{keepsection}\subsection{Exercise 5.2: Sierpinski triangle}

\begin{workedbox}{5.2.1 Continue the triangle pattern \ensuremath{\cdot} worked solution}
\Needspace{14\baselineskip}
\begin{enumerate}
\item Start at \textit{n} = 0 with one solid triangle. At every next step, split each black triangle into four equal smaller triangles and keep the three corner triangles black; the middle triangle stays blank.
\item \Needspace{11\baselineskip}
So the counts at the requested steps are:

\Needspace{10\baselineskip}
\begin{center}\small\begin{tabularx}{\linewidth}{>{\raggedright\arraybackslash}X>{\raggedright\arraybackslash}X>{\raggedright\arraybackslash}X}\toprule
\textbf{Step \textit{n}} & \textbf{Black triangles} & \textbf{Construction} \\ \midrule
0 & 1 & One solid starting triangle. \\
1 & 3 & Three black corner triangles. \\
2 & 9 & Each of the 3 black triangles produces 3 black corners. \\
3 & 27 & Each of the 9 black triangles produces 3 black corners. \\
\bottomrule\end{tabularx}\end{center}
\item For the drawing, repeat that replacement inside every black triangle at \textit{n} = 1 to obtain step 2, then repeat it inside all nine black triangles to obtain step 3. Previously empty central triangles remain empty.
\item \textbf{Check:} the number triples at each step, and the large central gaps remain blank rather than being recoloured.
\end{enumerate}
\end{workedbox}

\end{keepsection}
\begin{keepsection}\subsection{Exercise 5.2: Sierpinski triangle}

\begin{workedbox}{5.2.2 Write a recursive triangle counter \ensuremath{\cdot} worked solution}
\Needspace{13\baselineskip}
\begin{enumerate}
\item Let \texttt{black\_triangles(n)} return the number at step \textit{n}. The domain is a non-negative integer step, and step zero contains one black triangle.
\item \Needspace{5\baselineskip}
For \textit{n} \ensuremath{>} 0, each black triangle creates three black corner triangles, so the total is three times the previous step's count:

\begin{equation*}\mathrm{T}(0)\, =\, 1,\, \quad \, \mathrm{T}(\mathrm{n})\, =\, 3\, \cdot \, \mathrm{T}(\mathrm{n}\, -\, 1)\end{equation*}
\item \Needspace{3\baselineskip}
That recurrence gives a fully recursive implementation:
\begin{tcolorbox}[colback=Paperblue,colframe=Paperblue,boxrule=0pt,arc=0pt,left=7pt,right=7pt,top=5pt,bottom=5pt]\small\sloppy
\noindent\hspace*{0em}\texttt{def black\_triangles(n):}\par
\noindent\hspace*{1em}\texttt{if n == 0:}\par
\noindent\hspace*{3em}\texttt{return 1}\par
\noindent\hspace*{1em}\texttt{return 3 * black\_triangles(n - 1)}\par
\end{tcolorbox}
\item The base case returns directly. Each recursive call reduces \texttt{n}, and the returned smaller count is multiplied by three on the way back.
\item \textbf{Check:} \texttt{black\_triangles(0)} returns 1 and \texttt{black\_triangles(1)} returns 3, matching the first two drawings.
\end{enumerate}
\end{workedbox}

\end{keepsection}
\begin{keepsection}\subsection{Exercise 5.2: Sierpinski triangle}

\begin{workedbox}{5.2.3 Count the triangles at step 6 \ensuremath{\cdot} worked solution}
\Needspace{6\baselineskip}
\begin{enumerate}
\item \Needspace{13\baselineskip}
Continue multiplying the previous count by three:

\Needspace{10\baselineskip}
\begin{center}\small\begin{tabularx}{\linewidth}{>{\raggedright\arraybackslash}X>{\raggedright\arraybackslash}X>{\raggedright\arraybackslash}X}\toprule
\textbf{\textit{n}} & \textbf{Calculation} & \textbf{\textit{T}(\textit{n})} \\ \midrule
0 & Base case & 1 \\
1 & 3 \ensuremath{\times} 1 & 3 \\
2 & 3 \ensuremath{\times} 3 & 9 \\
3 & 3 \ensuremath{\times} 9 & 27 \\
4 & 3 \ensuremath{\times} 27 & 81 \\
5 & 3 \ensuremath{\times} 81 & 243 \\
6 & 3 \ensuremath{\times} 243 & 729 \\
\bottomrule\end{tabularx}\end{center}
\item Equivalently, \textit{T}(6) = 3\textsuperscript{6} = \textbf{729 black triangles}.
\item \textbf{Check:} six recursive steps multiply the starting count 1 by 3 six times, giving 729.
\end{enumerate}
\end{workedbox}

\end{keepsection}
\begin{keepsection}\subsection{Exercise 5.3: Recursion application}

\begin{workedbox}{5.3.1 Recursion and repetition \ensuremath{\cdot} worked solution}
\Needspace{6\baselineskip}
\begin{enumerate}
\item \textbf{Recursion} repeats work through a function calling itself on a smaller problem, whereas \textbf{repetition} repeats a block with a loop.
\item \textbf{Check:} both can perform work repeatedly; the distinction is how the next repetition is expressed.
\end{enumerate}
\end{workedbox}

\end{keepsection}
\begin{keepsection}\subsection{Exercise 5.3: Recursion application}

\begin{workedbox}{5.3.2 Write the recursive \texttt{dream} function \ensuremath{\cdot} worked solution}
\Needspace{22\baselineskip}
\begin{enumerate}
\item Assume level 0 means awake, \texttt{current\_level} is the level the dreamer is in now, and the helper's \texttt{level} argument names the destination level. Start with \texttt{dream(0, max\_level)}, where \texttt{max\_level} is non-negative.
\item The recursive case enters the next deeper level with \texttt{play\_lullaby(current\_level + 1)}, then calls \texttt{dream} for that level. The base case runs \texttt{do\_important\_job()} once at the maximum level.
\item When the deeper call returns, \texttt{kick(current\_level)} wakes the dreamer to this call's level. That line runs after the recursive call, so it executes on the way back up the call stack.
\item \Needspace{4\baselineskip}
Here is one complete version. It rejects a maximum below the current level, returns totems only from the outer awake-level call, and checks the totem after the last upward transition:
\begin{tcolorbox}[colback=Paperblue,colframe=Paperblue,boxrule=0pt,arc=0pt,left=7pt,right=7pt,top=5pt,bottom=5pt]\small\sloppy
\noindent\hspace*{0em}\texttt{def dream(current\_level, max\_level):}\par
\noindent\hspace*{1em}\texttt{if current\_level \ensuremath{<} 0 or max\_level \ensuremath{<} current\_level:}\par
\noindent\hspace*{3em}\texttt{raise ValueError("expected 0 \ensuremath{<}= current\_level \ensuremath{<}= max\_level")}\par
\noindent\hspace*{0em}\texttt{}\par
\noindent\hspace*{1em}\texttt{if current\_level == max\_level:}\par
\noindent\hspace*{3em}\texttt{do\_important\_job()}\par
\noindent\hspace*{1em}\texttt{else:}\par
\noindent\hspace*{3em}\texttt{next\_level = current\_level + 1}\par
\noindent\hspace*{3em}\texttt{play\_lullaby(next\_level)}\par
\noindent\hspace*{3em}\texttt{dream(next\_level, max\_level)}\par
\noindent\hspace*{3em}\texttt{kick(current\_level)}\par
\noindent\hspace*{0em}\texttt{}\par
\noindent\hspace*{1em}\texttt{if current\_level == 0:}\par
\noindent\hspace*{3em}\texttt{return check\_totem()}\par
\end{tcolorbox}
\item For \texttt{dream(0, 3)}, the calls go 0 \ensuremath{\to} 1 \ensuremath{\to} 2 \ensuremath{\to} 3. The job runs at 3; returns then call \texttt{kick(2)}, \texttt{kick(1)}, and \texttt{kick(0)}. Once awake at 0, \texttt{check\_totem()} runs.
\item \textbf{Check:} \texttt{play\_lullaby} runs once per downward move, \texttt{kick} once per upward move, the job runs once at the maximum, and the totem check runs only after the outermost call wakes.
\end{enumerate}
\end{workedbox}

\end{keepsection}

\vfill{\small Personal worked solutions based on APDR Blackboard Exercise 5, written in my own words; no exercise-sheet prose or figures are reproduced. The dream helpers are abstract functions supplied by the prompt. Not an official FAU publication; check course details against the current materials. Made by \textbf{Thanate Amornjiraporn (Focus)}.}

\end{document}
